\chapter{Chapter05 — Gotsman Linial Extremizer}
%%% generated by the blueprint pipeline from FABL.Chapter05.GotsmanLinialExtremizer
\section{Overview}

Declarations of the Lean module \texttt{FABL.Chapter05.GotsmanLinialExtremizer}.

\section{Declarations}

\begin{definition}[Gotsman–Linial starting index]
\lean{FABL.gotsmanLinialStart}
\label{FABL.gotsmanLinialStart}
\leanok
For natural numbers $n$ and $k$, define the starting index to be the natural-number
quotient of $n-k$ by $2$; here subtraction is truncated at zero.
\end{definition}

\begin{definition}[Gotsman–Linial root locations]
\lean{FABL.gotsmanLinialRoot}
\label{FABL.gotsmanLinialRoot}
\leanok
\uses{FABL.gotsmanLinialStart}
For natural numbers $n$, $k$, and $j$, let $a$ be the natural-number quotient of $n-k$
by $2$. Define the corresponding real number by $2(a+j)-n+1$.
\end{definition}

\begin{definition}[Gotsman–Linial polynomial]
\lean{FABL.gotsmanLinialPolynomial}
\statementsource{Analysis-of-Boolean-Functions-by-Ryan-ODonnell}{pdf-ffbe9c7b3306-p133-b004}
\label{FABL.gotsmanLinialPolynomial}
\leanok
\uses{FABL.gotsmanLinialRoot}
For natural numbers $n$ and $k$, define a real polynomial by
$P_{n,k}(X)=\prod_{j=0}^{k-1}\bigl(X-r_{n,k}(j)\bigr)$, where $r_{n,k}(j)$ is the
Gotsman–Linial root associated with $n$, $k$, and $j$.
\end{definition}

\begin{definition}[Gotsman–Linial polynomial on the cube]
\lean{FABL.gotsmanLinialCubePolynomial}
\label{FABL.gotsmanLinialCubePolynomial}
\leanok
\uses{FABL.gotsmanLinialPolynomial, FABL.signValue}
For natural numbers $n$ and $k$, define a function on sign vectors $x\in\{-1,1\}^n$ by
evaluating the Gotsman–Linial polynomial with parameters $n$ and $k$ at the sum of the
coordinate signs of $x$.
\end{definition}

\begin{definition}[Gotsman–Linial extremizer]
\lean{FABL.gotsmanLinialExtremizer}
\statementsource{Analysis-of-Boolean-Functions-by-Ryan-ODonnell}{pdf-ffbe9c7b3306-p133-b004}
\label{FABL.gotsmanLinialExtremizer}
\leanok
\uses{FABL.BooleanFunction, FABL.gotsmanLinialCubePolynomial, FABL.thresholdSign}
For natural numbers $n$ and $k$, the Gotsman–Linial extremizer is the Boolean function
on the $n$-dimensional Boolean cube that sends each input $x$ to the sign of the
Gotsman–Linial cube polynomial evaluated at $x$.
\end{definition}

\begin{theorem}[Symmetry of the extremizing function]
\lean{FABL.gotsmanLinialExtremizer_isSymmetric}
\proofsource{Analysis-of-Boolean-Functions-by-Ryan-ODonnell}{pdf-ffbe9c7b3306-p133-b004}
\label{FABL.gotsmanLinialExtremizer_isSymmetric}
\leanok
\uses{FABL.IsSymmetric, FABL.gotsmanLinialCubePolynomial, FABL.gotsmanLinialExtremizer, FABL.signValue, FABL.thresholdSign}
\difficulty{3}
For all natural numbers $n,k$, the function on $\{\pm1\}^n$ that assigns to $x$ the sign
of the polynomial $\prod_{j=0}^{k-1}\bigl(S(x)-[2(a+j)-n+1]\bigr)$ is invariant under
every permutation of the coordinates, where $S(x)=\sum_{i=1}^n x_i$ and
$a=\lfloor(n-k)/2\rfloor$; the sign is $1$ for a nonnegative value and $-1$ otherwise.
\end{theorem}

\begin{theorem}[Degree of the Gotsman–Linial polynomial]
\lean{FABL.natDegree_gotsmanLinialPolynomial}
\label{FABL.natDegree_gotsmanLinialPolynomial}
\leanok
\uses{FABL.gotsmanLinialPolynomial}
\difficulty{2}
For all natural numbers $n$ and $k$, the polynomial formed as the product of the $k$
linear factors associated with the specified Gotsman–Linial roots has natural degree
$k$.
\end{theorem}

\begin{theorem}[Multiplication by a coordinate sign]
\lean{FABL.fourierCoeff_mul_signValue}
\proofsource{Analysis-of-Boolean-Functions-by-Ryan-ODonnell}{
  pdf-ffbe9c7b3306-p024-b005,
  pdf-ffbe9c7b3306-p024-b006
}
\label{FABL.fourierCoeff_mul_signValue}
\leanok
\uses{FABL.fourierCoeff, FABL.monomial, FABL.monomial_mul_monomial, FABL.signValue}
\difficulty{3}
Let $q$ be any real-valued function on the $n$-dimensional sign cube, let $i$ be a
coordinate, and let $S$ be a set of coordinates. The Fourier coefficient at $S$ of the
function $x\mapsto q(x)x_i$ equals the Fourier coefficient of $q$ at the symmetric
difference of $S$ and $\{i\}$.
\end{theorem}

\begin{theorem}[Cardinality under symmetric difference]
\lean{FABL.card_le_card_symmDiff_singleton_add_one}
\label{FABL.card_le_card_symmDiff_singleton_add_one}
\leanok
\difficulty{4}
Let $n$ be a natural number, let $S\subseteq\{0,\ldots,n-1\}$, and let
$i\in\{0,\ldots,n-1\}$. Then $|S|\leq |S\mathbin{\triangle}\{i\}|+1$.
\end{theorem}

\begin{theorem}[Degree under a coordinate sum factor]
\lean{FABL.fourierDegree_mul_coordinateSum_sub_le}
\label{FABL.fourierDegree_mul_coordinateSum_sub_le}
\leanok
\uses{FABL.card_le_card_symmDiff_singleton_add_one, FABL.fourierCoeff, FABL.fourierCoeff_mul_signValue, FABL.fourierDegree, FABL.fourierDegree_le_iff, FABL.monomial, FABL.signValue}
\difficulty{5}
Let $q$ be a real-valued function on the Boolean cube with Fourier degree at most $d$,
and let $r\\in\\mathbb R$. The function $x\\mapsto q(x)\\bigl(\\sum_{i=1}^n
x_i-r\\bigr)$ has Fourier degree at most $d+1$.
\end{theorem}

\begin{definition}[Gotsman–Linial product on the cube]
\lean{FABL.gotsmanLinialProductOnCube}
\label{FABL.gotsmanLinialProductOnCube}
\leanok
\uses{FABL.signValue}
For natural numbers $n,a,d$ and a sign vector $x\in\{-1,1\}^n$, define a real-valued
function by taking the product over $j=0,\ldots,d-1$ of
$\sum_{i=0}^{n-1}\operatorname{sgn}(x_i)-\bigl(2(a+j)-n+1\bigr)$, where
$\operatorname{sgn}(x_i)$ is the sign value of coordinate $i$.
\end{definition}

\begin{theorem}[Degree of a cube product polynomial]
\lean{FABL.fourierDegree_gotsmanLinialProductOnCube_le}
\label{FABL.fourierDegree_gotsmanLinialProductOnCube_le}
\leanok
\uses{FABL.expect_monomial, FABL.fourierCoeff, FABL.fourierDegree, FABL.fourierDegree_le_iff, FABL.fourierDegree_mul_coordinateSum_sub_le, FABL.gotsmanLinialProductOnCube, FABL.signValue}
\difficulty{4}
For natural numbers $n,a,d$, define a function on $\{-1,1\}^n$ by taking the product
over $j=0,\ldots,d-1$ of $\sum_{i=0}^{n-1}x_i-\bigl(2(a+j)-n+1\bigr)$. Its Fourier
degree is at most $d$.
\end{theorem}

\begin{theorem}[Cube polynomial product formula]
\lean{FABL.gotsmanLinialCubePolynomial_eq_product}
\label{FABL.gotsmanLinialCubePolynomial_eq_product}
\leanok
\uses{FABL.gotsmanLinialCubePolynomial, FABL.gotsmanLinialPolynomial, FABL.gotsmanLinialProductOnCube, FABL.gotsmanLinialRoot, FABL.gotsmanLinialStart}
\difficulty{3}
For natural numbers $n,k$ and a sign vector $x\in\{\pm1\}^n$, let $S(x)=\sum_{i=1}^n
x_i$ and $a=\lfloor(n-k)/2\rfloor$. The polynomial obtained by evaluating
$\prod_{j=0}^{k-1}\bigl(X-[2(a+j)-n+1]\bigr)$ at $S(x)$ equals
$\prod_{j=0}^{k-1}\bigl(S(x)-[2(a+j)-n+1]\bigr)$.
\end{theorem}

\begin{theorem}[Gotsman linial extremizer is polynomial threshold]
\lean{FABL.gotsmanLinialExtremizer_isPolynomialThreshold}
\statementsource{Analysis-of-Boolean-Functions-by-Ryan-ODonnell}{
  pdf-ffbe9c7b3306-p115-b005,
  pdf-ffbe9c7b3306-p133-b004
}
\label{FABL.gotsmanLinialExtremizer_isPolynomialThreshold}
\leanok
\uses{FABL.IsPolynomialThreshold, FABL.fourierDegree_gotsmanLinialProductOnCube_le, FABL.gotsmanLinialCubePolynomial, FABL.gotsmanLinialCubePolynomial_eq_product, FABL.gotsmanLinialExtremizer}
\difficulty{3}
% dataset-pipeline: stub -- prose pending build_informal_statements.py
Statement pending regeneration.
\end{theorem}

\begin{theorem}[Nonvanishing of each polynomial factor]
\lean{FABL.gotsmanLinial_factor_ne_zero}
\label{FABL.gotsmanLinial_factor_ne_zero}
\leanok
\uses{FABL.gotsmanLinialRoot, FABL.gotsmanLinialStart, FABL.positiveCoordinateCount, FABL.signValue, FABL.sum_signValue_eq_two_mul_positiveCoordinateCount_sub}
\difficulty{4}
Let $n,k$ be natural numbers, let $x\in\{\pm1\}^n$, and let $j$ be a natural number with
$j<k$. Then $\sum_{i=1}^n x_i-[2(a+j)-n+1]\ne0$, where $a=\lfloor(n-k)/2\rfloor$.
\end{theorem}

\begin{theorem}[Nonvanishing of the cube polynomial]
\lean{FABL.gotsmanLinialCubePolynomial_ne_zero}
\label{FABL.gotsmanLinialCubePolynomial_ne_zero}
\leanok
\uses{FABL.gotsmanLinialCubePolynomial, FABL.gotsmanLinialCubePolynomial_eq_product, FABL.gotsmanLinialProductOnCube, FABL.gotsmanLinialRoot, FABL.gotsmanLinial_factor_ne_zero}
\difficulty{3}
For every pair of natural numbers $n,k$ and every sign vector $x\in\{\pm1\}^n$, the
value at $S(x)=\sum_{i=1}^n x_i$ of the polynomial
$\prod_{j=0}^{k-1}\bigl(X-[2(a+j)-n+1]\bigr)$ is nonzero, where
$a=\lfloor(n-k)/2\rfloor$.
\end{theorem}

\begin{theorem}[Positivity of the signed central product]
\lean{FABL.centralProduct_signed_pos}
\label{FABL.centralProduct_signed_pos}
\leanok
\difficulty{5}
For all nonnegative integers $j$ and $d$, $\displaystyle
(-1)^d\prod_{t=0}^{j+d-1}(2j-2t-1)>0$, where an empty product is $1$.
\end{theorem}

\begin{theorem}[Threshold sign of a central product]
\lean{FABL.thresholdSign_centralProduct}
\label{FABL.thresholdSign_centralProduct}
\leanok
\uses{FABL.centralProduct_signed_pos, FABL.thresholdSign, FABL.thresholdSign_of_neg, FABL.thresholdSign_of_nonneg}
\difficulty{4}
For all natural numbers $j$ and $d$, the threshold sign of
$\\prod_{t=0}^{j+d-1}(2j-2t-1)$ is $1$ when $d$ is even and $-1$ when $d$ is odd.
\end{theorem}

\begin{theorem}[Gotsman linial extremizer central layer]
\lean{FABL.gotsmanLinialExtremizer_centralLayer}
\statementsource{Analysis-of-Boolean-Functions-by-Ryan-ODonnell}{pdf-ffbe9c7b3306-p133-b004}
\label{FABL.gotsmanLinialExtremizer_centralLayer}
\leanok
\uses{FABL.canonicalCountInput, FABL.gotsmanLinialCubePolynomial_eq_product, FABL.gotsmanLinialExtremizer, FABL.gotsmanLinialProductOnCube, FABL.gotsmanLinialStart, FABL.positiveCoordinateCount, FABL.positiveCoordinateCount_canonicalCountInput, FABL.sum_signValue_eq_two_mul_positiveCoordinateCount_sub, FABL.thresholdSign, FABL.thresholdSign_centralProduct}
\difficulty{4}
% dataset-pipeline: stub -- prose pending build_informal_statements.py
Statement pending regeneration.
\end{theorem}

\begin{definition}[Gotsman–Linial layer product]
\lean{FABL.gotsmanLinialLayerValue}
\label{FABL.gotsmanLinialLayerValue}
\leanok
For natural numbers $a$, $k$, and $c$, define the real number obtained by multiplying
$2c-2(a+j)-1$ over all integers $j$ with $0 \le j < k$.
\end{definition}

\begin{theorem}[Nonvanishing of the layer product]
\lean{FABL.gotsmanLinialLayerValue_ne_zero}
\label{FABL.gotsmanLinialLayerValue_ne_zero}
\leanok
\uses{FABL.gotsmanLinialLayerValue}
\difficulty{3}
For all natural numbers $a,k,c$, the product $\prod_{j=0}^{k-1}\bigl(2c-2(a+j)-1\bigr)$
is nonzero.
\end{theorem}

\begin{theorem}[Layer product count recurrence]
\lean{FABL.gotsmanLinialLayerValue_succ_count}
\label{FABL.gotsmanLinialLayerValue_succ_count}
\leanok
\uses{FABL.gotsmanLinialLayerValue}
\difficulty{3}
For all natural numbers $a,s,c$, the layer product with parameters $(a,s+1,c+1)$ equals
$\bigl(2c-2a+1\bigr)$ times the layer product with parameters $(a,s,c)$; here the layer
product at $(a,k,c)$ is $\prod_{j=0}^{k-1}\bigl(2c-2(a+j)-1\bigr)$.
\end{theorem}

\begin{theorem}[Layer product degree recurrence]
\lean{FABL.gotsmanLinialLayerValue_succ_degree}
\label{FABL.gotsmanLinialLayerValue_succ_degree}
\leanok
\uses{FABL.gotsmanLinialLayerValue}
\difficulty{3}
For all natural numbers $a,s,c$, the layer product with parameters $(a,s+1,c)$ equals
the layer product with parameters $(a,s,c)$ multiplied by $2c-2a-2s-1$, where the layer
product at $(a,k,c)$ is $\prod_{j=0}^{k-1}\bigl(2c-2(a+j)-1\bigr)$.
\end{theorem}

\begin{theorem}[Layer values with different threshold signs]
\lean{FABL.thresholdSign_gotsmanLinialLayerValue_succ_ne_iff}
\label{FABL.thresholdSign_gotsmanLinialLayerValue_succ_ne_iff}
\leanok
\uses{FABL.gotsmanLinialLayerValue, FABL.gotsmanLinialLayerValue_ne_zero, FABL.gotsmanLinialLayerValue_succ_count, FABL.gotsmanLinialLayerValue_succ_degree, FABL.thresholdSign, FABL.thresholdSign_of_neg, FABL.thresholdSign_of_nonneg}
\difficulty{6}
For natural numbers $a,s,c$, define $V(a,k,c)=\\prod_{j=0}^{k-1}(2c-2(a+j)-1)$. The
threshold signs of $V(a,s+1,c+1)$ and $V(a,s+1,c)$ differ if and only if $a\\le
c<a+s+1$.
\end{theorem}

\begin{theorem}[Sign sum on a dimension edge]
\lean{FABL.sum_signValue_dimensionEdge_plus}
\label{FABL.sum_signValue_dimensionEdge_plus}
\leanok
\uses{FABL.DimensionEdge, FABL.dimensionEdgeRemoveEquiv, FABL.signValue}
\difficulty{3}
Let $m$ be a nonnegative integer, choose a coordinate $i$ among $m+1$ coordinates, and
let $x\in\{-1,1\}^{m+1}$ satisfy $x_i=1$. The sum of the signs of all coordinates of $x$
is $1$ plus the sum of the signs in the $m$-coordinate vector obtained by deleting
coordinate $i$.
\end{theorem}

\begin{theorem}[Sign sum on a dimension edge after flipping]
\lean{FABL.sum_signValue_dimensionEdge_minus}
\label{FABL.sum_signValue_dimensionEdge_minus}
\leanok
\uses{FABL.DimensionEdge, FABL.dimensionEdgeRemoveEquiv, FABL.flipCoordinate, FABL.setCoordinate, FABL.signValue}
\difficulty{3}
Let $m$ be a nonnegative integer, choose a coordinate $i$ among $m+1$ coordinates, and
let $x\in\{-1,1\}^{m+1}$ satisfy $x_i=1$. For a sign $s\in\{-1,1\}$ write its real value
as $s$. The sum of the signs of the coordinates of $x$ after flipping coordinate $i$ is
$-1$ plus the sum of the signs in the $m$-coordinate vector obtained by deleting
coordinate $i$ from $x$.
\end{theorem}

\begin{theorem}[Boundary edges of the Gotsman–Linial extremizer]
\lean{FABL.isBoundaryDimensionEdge_gotsmanLinialExtremizer_succ_iff}
\label{FABL.isBoundaryDimensionEdge_gotsmanLinialExtremizer_succ_iff}
\leanok
\uses{FABL.DimensionEdge, FABL.IsBoundaryDimensionEdge, FABL.dimensionEdgeRemoveEquiv, FABL.flipCoordinate, FABL.gotsmanLinialCubePolynomial, FABL.gotsmanLinialCubePolynomial_eq_product, FABL.gotsmanLinialExtremizer, FABL.gotsmanLinialLayerValue, FABL.gotsmanLinialProductOnCube, FABL.gotsmanLinialStart, FABL.positiveCoordinateCount, FABL.signValue, FABL.sum_signValue_dimensionEdge_minus, FABL.sum_signValue_dimensionEdge_plus, FABL.sum_signValue_eq_two_mul_positiveCoordinateCount_sub, FABL.thresholdSign, FABL.thresholdSign_gotsmanLinialLayerValue_succ_ne_iff}
\difficulty{5}
Let $m$ and $s$ be natural numbers, choose a coordinate $i$ among $m+1$ coordinates, and
let $e$ be a cube point whose $i$th coordinate is $1$. The Gotsman–Linial extremizer
with parameters $m+1$ and $s+1$ changes value across the edge at $e$ in coordinate $i$
if and only if, after removing coordinate $i$, the number of coordinates equal to $1$ is
$(m-s)_+/2$ rounded down, plus some integer $j$ with $0\le j<s+1$; here
$(m-s)_+=\max(m-s,0)$.
\end{theorem}

\begin{theorem}[Probability of a coordinate-count interval]
\lean{FABL.uniformProbability_exists_positiveCoordinateCount_eq_add}
\label{FABL.uniformProbability_exists_positiveCoordinateCount_eq_add}
\leanok
\uses{FABL.positiveCoordinateCount, FABL.uniformProbability, FABL.uniformProbability_positiveCoordinateCount_eq}
\difficulty{5}
For nonnegative integers $m,a,k$, the uniform probability that a sign vector
$y\\in\\{-1,1\\}^m$ has between $a$ and $a+k-1$ positive coordinates, inclusive, is
$\\displaystyle\\sum_{j=0}^{k-1}\\frac{\\binom{m}{a+j}}{2^m}$.
\end{theorem}

\begin{theorem}[Influence of the Gotsman–Linial extremizer]
\lean{FABL.booleanInfluence_gotsmanLinialExtremizer}
\label{FABL.booleanInfluence_gotsmanLinialExtremizer}
\leanok
\uses{FABL.IsBoundaryDimensionEdge, FABL.IsPivotal, FABL.booleanInfluence, FABL.booleanInfluence_eq_dimensionEdgeBoundaryFraction, FABL.dimensionEdgeBoundaryFraction, FABL.dimensionEdgeRemoveEquiv, FABL.gotsmanLinialCubePolynomial, FABL.gotsmanLinialExtremizer, FABL.gotsmanLinialPolynomial, FABL.gotsmanLinialStart, FABL.isBoundaryDimensionEdge_gotsmanLinialExtremizer_succ_iff, FABL.positiveCoordinateCount, FABL.uniformProbability, FABL.uniformProbability_exists_positiveCoordinateCount_eq_add}
\difficulty{4}
For nonnegative integers $m$ and $k$, and any coordinate $i$ of the $(m+1)$-dimensional
Boolean cube, the influence of coordinate $i$ for the Gotsman–Linial extremizer with
parameters $m+1$ and $k$ is
$$\\sum_{j=0}^{k-1}\\frac{\\binom{m}{\\lfloor((m+1)-k)/2\\rfloor+j}}{2^m}.$$
\end{theorem}

\begin{theorem}[Total influence gotsman linial extremizer]
\lean{FABL.totalInfluence_gotsmanLinialExtremizer}
\statementsource{Analysis-of-Boolean-Functions-by-Ryan-ODonnell}{pdf-ffbe9c7b3306-p133-b004}
\label{FABL.totalInfluence_gotsmanLinialExtremizer}
\leanok
\uses{FABL.BooleanFunction.toReal, FABL.booleanInfluence_eq_influence_toReal, FABL.booleanInfluence_gotsmanLinialExtremizer, FABL.gotsmanLinialExtremizer, FABL.gotsmanLinialStart, FABL.totalInfluence}
\difficulty{3}
% dataset-pipeline: stub -- prose pending build_informal_statements.py
Statement pending regeneration.
\end{theorem}

\begin{definition}[Shifted binomial probability]
\lean{FABL.shiftedBinomialProbability}
\label{FABL.shiftedBinomialProbability}
\leanok
For natural numbers $m$, $c$, and $j$, define the shifted binomial probability as
$\binom{2m+c}{m+j}/2^{2m+c}$, where the binomial coefficient is zero if $m+j>2m+c$.
\end{definition}

\begin{theorem}[Initial shifted binomial probability]
\lean{FABL.shiftedBinomialProbability_zero}
\label{FABL.shiftedBinomialProbability_zero}
\leanok
\uses{FABL.oddMajorityInfluence, FABL.shiftedBinomialProbability}
\difficulty{1}
For every natural number $m$, the shifted binomial probability $\binom{2m}{m}/2^{2m}$
equals the odd majority influence $\binom{2m}{m}/2^{2m}$.
\end{theorem}

\begin{theorem}[Shifted binomial probability recurrence]
\lean{FABL.shiftedBinomialProbability_succ_same}
\label{FABL.shiftedBinomialProbability_succ_same}
\leanok
\uses{FABL.shiftedBinomialProbability}
\difficulty{4}
For natural numbers $m,c,j$ with $j\le c$, let $P(m,c,j)=\binom{2m+c}{m+j}/2^{2m+c}$.
Then $P(m,c+1,j)=P(m,c,j)\,\frac{2m+c+1}{2(m+c+1-j)}$.
\end{theorem}

\begin{theorem}[Diagonal shifted binomial recurrence]
\lean{FABL.shiftedBinomialProbability_succ_diagonal}
\label{FABL.shiftedBinomialProbability_succ_diagonal}
\leanok
\uses{FABL.shiftedBinomialProbability}
\difficulty{4}
For natural numbers $m,c$, let $P(m,c,j)=\binom{2m+c}{m+j}/2^{2m+c}$. Then
$P(m,c+1,c+1)=P(m,c,c)\,\frac{2m+c+1}{2(m+c+1)}$.
\end{theorem}

\begin{theorem}[Shifted binomial factor limit]
\lean{FABL.tendsto_shiftedBinomialSameFactor}
\label{FABL.tendsto_shiftedBinomialSameFactor}
\leanok
\difficulty{3}
For fixed natural numbers $c$ and $j$ with $j\le c$, the quantity
$(2m+c+1)/(2(m+c+1-j))$ tends to $1$ as the natural number $m$ tends to infinity.
\end{theorem}

\begin{theorem}[Limit of a shifted diagonal factor]
\lean{FABL.tendsto_shiftedBinomialDiagonalFactor}
\label{FABL.tendsto_shiftedBinomialDiagonalFactor}
\leanok
\difficulty{3}
For each natural number $c$, the real sequence indexed by natural numbers $m$ and given
by $\frac{2m+c+1}{2(m+c+1)}$ converges to $1$ as $m$ tends to infinity.
\end{theorem}

\begin{theorem}[Shifted binomial probability asymptotics]
\lean{FABL.tendsto_shiftedBinomialProbability_div_oddMajorityInfluence}
\label{FABL.tendsto_shiftedBinomialProbability_div_oddMajorityInfluence}
\leanok
\uses{FABL.oddMajorityInfluence, FABL.oddMajorityInfluence_pos, FABL.shiftedBinomialProbability, FABL.shiftedBinomialProbability_succ_diagonal, FABL.shiftedBinomialProbability_succ_same, FABL.shiftedBinomialProbability_zero, FABL.tendsto_shiftedBinomialDiagonalFactor, FABL.tendsto_shiftedBinomialSameFactor}
\difficulty{5}
For fixed natural numbers $c$ and $j$ with $j\le c$, let
$P_m=\binom{2m+c}{m+j}/2^{2m+c}$ and $I_m=\binom{2m}{m}/2^{2m}$. As $m\to\infty$, the
ratio $P_m/I_m$ tends to $1$.
\end{theorem}

\begin{theorem}[Tendsto sqrt shift mul odd majority influence]
\lean{FABL.tendsto_sqrt_shift_mul_oddMajorityInfluence}
\statementsource{Analysis-of-Boolean-Functions-by-Ryan-ODonnell}{
  pdf-ffbe9c7b3306-p062-b012,
  pdf-ffbe9c7b3306-p047-b001
}
\label{FABL.tendsto_sqrt_shift_mul_oddMajorityInfluence}
\leanok
\uses{FABL.oddMajorityInfluence, FABL.oddMajorityInfluenceMain, FABL.tendsto_oddMajorityInfluence_div_main}
\difficulty{6}
% dataset-pipeline: stub -- prose pending build_informal_statements.py
Statement pending regeneration.
\end{theorem}

\begin{theorem}[Shifted binomial sum limit]
\lean{FABL.tendsto_sum_shiftedBinomialProbability_div_oddMajorityInfluence}
\label{FABL.tendsto_sum_shiftedBinomialProbability_div_oddMajorityInfluence}
\leanok
\uses{FABL.oddMajorityInfluence, FABL.shiftedBinomialProbability, FABL.tendsto_shiftedBinomialProbability_div_oddMajorityInfluence}
\difficulty{3}
For natural numbers $c$ and $k$ with $k\le c+1$, and each natural number $m$, set
$p_{m,c,j}=\binom{2m+c}{m+j}/2^{2m+c}$ and $a_m=\binom{2m}{m}/2^{2m}$. As $m$ tends to
infinity, $\sum_{j=0}^{k-1} p_{m,c,j}/a_m$ tends to $k$.
\end{theorem}

\begin{definition}[Normalized Gotsman–Linial total influence]
\lean{FABL.normalizedGotsmanLinialTotalInfluence}
\label{FABL.normalizedGotsmanLinialTotalInfluence}
\leanok
\uses{FABL.BooleanFunction.toReal, FABL.gotsmanLinialExtremizer, FABL.totalInfluence}
For natural numbers $k$ and $r$, define the normalized total influence of the
Gotsman–Linial extremizer on $k+r$ coordinates with parameter $k$ as its total influence
divided by the square root of $k+r$.
\end{definition}

\begin{theorem}[Even-parameter influence identity]
\lean{FABL.normalizedGotsmanLinialTotalInfluence_even}
\label{FABL.normalizedGotsmanLinialTotalInfluence_even}
\leanok
\uses{FABL.gotsmanLinialStart, FABL.normalizedGotsmanLinialTotalInfluence, FABL.oddMajorityInfluence, FABL.oddMajorityInfluence_pos, FABL.shiftedBinomialProbability, FABL.totalInfluence_gotsmanLinialExtremizer}
\difficulty{4}
Let $k$ and $m$ be natural numbers with $k>0$. For $M_m=\binom{2m}{m}/2^{2m}$ and
$Q_{m,c,j}=\binom{2m+c}{m+j}/2^{2m+c}$, the normalized total influence of the
Gotsman–Linial extremizer on the $(k+2m)$-dimensional Boolean cube satisfies
\[\frac{\operatorname{Inf}_{\mathrm{tot}}(G_{k+2m,k})}{\sqrt{k+2m}}=\sqrt{k+2m}\,M_m\sum_{j=0}^{k-1}\frac{Q_{m,k-1,j}}{M_m}.\]
\end{theorem}

\begin{theorem}[Odd-parameter influence identity]
\lean{FABL.normalizedGotsmanLinialTotalInfluence_odd}
\label{FABL.normalizedGotsmanLinialTotalInfluence_odd}
\leanok
\uses{FABL.gotsmanLinialStart, FABL.normalizedGotsmanLinialTotalInfluence, FABL.oddMajorityInfluence, FABL.oddMajorityInfluence_pos, FABL.shiftedBinomialProbability, FABL.totalInfluence_gotsmanLinialExtremizer}
\difficulty{4}
Let $k$ and $m$ be natural numbers with $k>0$. For $M_m=\binom{2m}{m}/2^{2m}$ and
$Q_{m,c,j}=\binom{2m+c}{m+j}/2^{2m+c}$, the normalized total influence of the
Gotsman–Linial extremizer on the $(k+2m+1)$-dimensional Boolean cube satisfies
\[\frac{\operatorname{Inf}_{\mathrm{tot}}(G_{k+2m+1,k})}{\sqrt{k+2m+1}}=\sqrt{k+2m+1}\,M_m\sum_{j=0}^{k-1}\frac{Q_{m,k,j}}{M_m}.\]
\end{theorem}

\begin{theorem}[Even subsequence limit of normalized influence]
\lean{FABL.tendsto_normalizedGotsmanLinialTotalInfluence_even}
\label{FABL.tendsto_normalizedGotsmanLinialTotalInfluence_even}
\leanok
\uses{FABL.normalizedGotsmanLinialTotalInfluence, FABL.normalizedGotsmanLinialTotalInfluence_even, FABL.oddMajorityInfluence, FABL.tendsto_sqrt_shift_mul_oddMajorityInfluence, FABL.tendsto_sum_shiftedBinomialProbability_div_oddMajorityInfluence}
\difficulty{4}
For every integer $k>0$, as $m$ tends to infinity through the nonnegative integers, the
total influence of the degree-$k$ Gotsman–Linial extremizer on the $(k+2m)$-dimensional
Boolean cube, divided by $\sqrt{k+2m}$, tends to $k\sqrt{2/\pi}$.
\end{theorem}

\begin{theorem}[Odd subsequence limit of normalized influence]
\lean{FABL.tendsto_normalizedGotsmanLinialTotalInfluence_odd}
\label{FABL.tendsto_normalizedGotsmanLinialTotalInfluence_odd}
\leanok
\uses{FABL.normalizedGotsmanLinialTotalInfluence, FABL.normalizedGotsmanLinialTotalInfluence_odd, FABL.oddMajorityInfluence, FABL.tendsto_sqrt_shift_mul_oddMajorityInfluence, FABL.tendsto_sum_shiftedBinomialProbability_div_oddMajorityInfluence}
\difficulty{4}
For every integer $k>0$, as $m$ tends to infinity through the nonnegative integers, the
total influence of the degree-$k$ Gotsman–Linial extremizer on the
$(k+2m+1)$-dimensional Boolean cube, divided by $\sqrt{k+2m+1}$, tends to
$k\sqrt{2/\pi}$.
\end{theorem}

\begin{theorem}[Convergence from even and odd subsequences]
\lean{FABL.tendsto_nat_of_even_odd}
\label{FABL.tendsto_nat_of_even_odd}
\leanok
\difficulty{4}
Let $f:\bbn\to\bbr$ and $L\in\bbr$. If both subsequences $f(2m)$ and $f(2m+1)$ tend to
$L$ as $m\to\infty$, then $f(n)$ tends to $L$ as $n\to\infty$.
\end{theorem}

\begin{theorem}[Asymptotic normalized total influence]
\lean{FABL.tendsto_normalizedGotsmanLinialTotalInfluence}
\label{FABL.tendsto_normalizedGotsmanLinialTotalInfluence}
\leanok
\uses{FABL.normalizedGotsmanLinialTotalInfluence, FABL.tendsto_nat_of_even_odd, FABL.tendsto_normalizedGotsmanLinialTotalInfluence_even, FABL.tendsto_normalizedGotsmanLinialTotalInfluence_odd, FABL.totalInfluence_gotsmanLinialExtremizer}
\difficulty{4}
For each fixed integer $k\ge 0$, as the integer $r$ tends to infinity, the total
influence of the degree-$k$ Gotsman–Linial extremizer on the $(k+r)$-dimensional Boolean
cube, divided by $\sqrt{k+r}$, tends to $k\sqrt{2/\pi}$.
\end{theorem}
